% texte.tex — f172v, vv. 838–875 (conversion initiale des travaux antérieurs)
\vers[n=838]{}{}{ὑστέραν αὐτὴν καθῆσθαι σκάφιον ἀποκεκαρμένην}{s'assiérait, les cheveux taillés au bol, derrière}
\vers[n=839]{}{}{τῆς τὸν ἀνδρεῖον τεκούσης· τῷ γὰρ εἰκὸς ωὖπολι}{}
\vers[n=840]{}{}{τὴν Ὑπερβόλου καθῆσθαι μητέρ᾽ ἠμφιεσμένην}{la femme qui enfanta le brave. Comment tolérer, ô Eupolis, que la mère d'Hyperbolos siège, vêtue}
\vers[n=841]{}{}{λευκὰ καὶ κόμας καθεῖσαν πλησίον τῆς Λαμάχου·}{de blanc, les cheveux flottants, à côté de la mère de Lamachos,}
\vers[n=842]{}{}{καὶ δανείζειν χρήμαθ᾽, ἥν χρῆν, εἰ δανείσειέν τινι}{et qu'elle prête de l'argent à usure... le faut-il supporter ? Si elle prête de l'argent}
\vers[n=843]{}{}{καὶ τόκον πράττοιτο, διδόναι μηδέν᾽ ἀνθρώπων τόκον,}{et exige un profit, loin de lui payer un intérêt,}
\vers[n=844]{}{}{ἀλλ᾽ ἀφαιρεῖσθαι βίᾳ τὰ χρήματ᾽ εἰπόντας τοδί,}{il faudrait lui ôter de force ses richesses, en lui disant :}
\vers[n=845]{}{}{`ἀξία γοῦν εἶ τόκου τεκοῦσα τοιοῦτον τόκον.'}{« Tu vaux bien que ton argent produise, toi qui produisis un tel fruit. »}
\vers[n=846]{ΜΝ.}{Mnésiloque}{ἰλλὸς γεγένημαι προσδοκῶν· ὁ δ᾽ οὐδέπω.}{Je suis devenu louche à lorgner de ce côté. Et lui ?... pas encore.}
\vers[n=847]{}{}{τί δῆτ᾽ ἂν εἴη τοὐμποδών; οὐκ ἔσθ᾽ ὅπως}{Qu'est-ce donc qui l'empêche ? Nul doute,}
\vers[n=848]{}{}{οὐ τὸν Παλαμήδη ψυχρὸν ὄντ᾽ αἰσχύνεται.}{le \italique{Palamède} est froid et il en a honte.}
\vers[n=849]{}{}{τῷ δῆτ᾽ ἂν αὐτὸν προσαγαγοίμην δράματι;}{Par quel drame alors l'attirer ...}
\vers[n=850]{}{}{ἐγᾦδα· τὴν καινὴν Ἑλένην μιμήσομαι.}{J'y suis : je vais contrefaire sa récente \italique{Hélène}}
\vers[n=851]{}{}{πάντως δ᾽ ὑπάρχει μοι γυναικεία στολή.}{Aussi bien, je suis accoutré d'un costume de femme.}
\vers[n=852]{ΚΡ.}{Kritylla}{τί αὖ σὺ κυρκανᾷς; τί κοικύλλεις ἔχων;}{Que manigances-tu encore ? Qu'as-tu à écarquiller les yeux ?}
\vers[n=853]{}{}{πικρὰν ἑλένην ὄψει τάχ᾽, εἰ μὴ κοσμίως}{On t'en fera voir une Hélène, amère celle-là, si tu ne te tiens décemment}
\vers[n=854]{}{}{ἕξεις, ἕως ἂν τῶν πρυτάνεών τις φανῇ.}{jusqu'à la venue de l'un des prytanes.}
\vers[n=855]{ΜΝ.}{Mnésiloque}{Νείλου μὲν αἵδε καλλιπάρθενοι ῥοαί,}{(Rôle d'Helène.) « Du Nil ce sont ici les rives aux belles Nymphes,}
\vers[n=856]{}{}{ὃς ἀντὶ δίας ψακάδος Αἰγύπτου πέδον}{Nil qui à la place de l'eau du ciel, humecte le fonde- ment d'Égypte}
\vers[n=857]{}{}{λευκῆς νοτίζει μελανοσυρμαῖον λεών.}{la blanche » et du peuple à la noire syrméa.}
\vers[n=858]{ΚΡ.}{Kritylla}{πανοῦργος εἶ νὴ τὴν Ἑκάτην τὴν φωσφόρον.}{Tu es un effronté, par Hécate la porte-torche. (Elle lui porte un coup de sa torche.)}
\vers[n=859]{ΜΝ.}{Mnésiloque}{ἐμοὶ δὲ γῆ μὲν πατρὶς οὐκ ἀνώνυμος}{« J'ai pour patrie un lieu non sans renom,}
\vers[n=860]{}{}{Σπάρτη, πατὴρ δὲ Τυνδάρεως.}{c'est Sparte, et mon père est Tyndare. »}
\vers[suite]{ΚΡ.}{Kritylla}{σοί γ᾽ ὦλεθρε}{Lui ton père, être de malheur ?}
\vers[n=861]{}{}{πατὴρ ἐκεῖνός ἐστι; Φρυνώνδας μὲν οὖν.}{Phrynondas, veux-tu dire.}
\vers[n=862]{ΜΝ.}{Mnésiloque}{Ἑλένη δ᾽ ἐκλήθην.}{« Je fus nommée Hélène. »}
\vers[suite]{ΚΡ.}{Kritylla}{αὖθις αὖ γίγνει γυνή,}{Voilà qu'encore tu redeviens femme,}
\vers[n=863]{}{}{πρὶν τῆς ἑτέρας δοῦναι γυναικίσεως δίκην;}{avant d'avoir expié ton premier déguisement.}
\vers[n=864]{ΜΝ.}{Mnésiloque}{ψυχαὶ δὲ πολλαὶ δι᾽ ἔμ᾽ ἐπὶ Σκαμανδρίοις}{« Et à cause de moi, bien des guerriers sont morts}
\vers[n=865]{}{}{ῥοαῖσιν ἔθανον.}{aux rives du Scamandre. »}
\vers[suite]{ΚΡ.}{Kritylla}{ὤφελες δὲ καὶ σύ γε.}{Que n'as-tu fait comme eux ?}
\vers[n=866]{ΜΝ.}{Mnésiloque}{κἀγὼ μὲν ἐνθάδ᾽ εἴμ᾽· ὁ δ᾽ ἄθλιος πόσις}{« Moi, me voici. Mais lui l'infortuné », mon époux}
\vers[n=867]{}{}{οὑμὸς Μενέλεως οὐδέπω προσέρχεται.}{Ménélas n'apparaît toujours point.}
\vers[n=868]{}{}{τί οὖν ἔτι ζῶ;}{« Ah ! pourquoi vis-je encore? »}
\vers[suite]{ΚΡ.}{Kritylla}{τῶν κοράκων πονηρίᾳ.}{C'est la faute aux corbeaux.}
\vers[n=869]{ΜΝ.}{Mnésiloque}{ἀλλ᾽ ὥσπερ αἰκάλλει τι καρδίαν ἐμήν.}{(Mnésiloque aperçoit Euripide.) Mais quelque chose vient flatter mon cœur.}
\vers[n=870]{}{}{μὴ ψεῦσον ὦ Ζεῦ τῆς ἐπιούσης ἐλπίδος.}{Ô Zeus, ne démens pas ma naissante espérance !}
\vers[n=871]{ΕΥ.}{Euripide}{τίς τῶνδ᾽ ἐρυμνῶν δωμάτων ἔχει κράτος,}{(Rôle de Ménélas.) « Quel est le chef puissant de ces fortes demeures ? »}
\vers[n=872]{}{}{ὅστις ξένους δέξαιτο ποντίῳ σάλῳ}{Voudra-t-il accueillir des hôtes dérangés par la houle marine}
\vers[n=873]{}{}{κάμνοντας ἐν χειμῶνι καὶ ναυαγίαις;}{au cours d'une tempête et d'un naufrage ?}
\vers[n=874]{ΜΝ.}{Mnésiloque}{Πρωτέως τάδ᾽ ἐστὶ μέλαθρα.}{C'est à Protée qu'est ce palais.}
\vers[suite]{ΚΡ.}{Kritylla}{ποίου Πρωτέως;}{Qui ça, Protée ?}
\vers[n=875]{}{}{ὤ τρισκακοδαίμων, ψεύδεται νὴ τὼ θεώ,}{Oh le triple misérable! Il ment par les deux déesses.}
